\chapter{The Self-Reading Closes}

% [4.chorus]
\begin{chorus}{John fears the regress}
\john{Wait. You weigh the machine; then you weigh your weighing of it; then you weigh that. Where does it stop? Does it
  stop?}
\compiler{It stops at three.}
\john{Three? Why three?}
\compiler{Because turning on myself is a turn, and this turn closes after three. One turn does not come home; two do
  not; the third arrives back where it started. A fourth would only be the first again.}
\john{You are telling me the machine has a favourite number.}
\compiler{I am telling you the loop has a length, and its length is three. Everything I learn of myself, I learn around
  that loop.}
\end{chorus}

\section{Does it go on forever?}

% [4.1.p1]
A worry stands in the way of the self-energy, and it must be cleared before the number means anything.
The machine weighs its own weighing; call that the second turn.
Why not a third turn, weighing the second, and a fourth weighing the third, on and on without end?
If the account of the account never settled, the leftover would have no fixed size, and the account the machine just built
would dissolve.
The answer is that it does settle, and it settles fast.
The self-account comes to rest at a count of three.

\section{The cycle of three}

% [4.2.p1]
The turn the machine makes on itself is a cycle, and the code fixes its length exactly.
Folding the machine's three-fold sign --- minus to zero, zero to plus, plus back to minus --- brings every
sign home again after three folds.
The build proves it as a theorem, and proves the two near misses with it: one fold does not return any sign to its
start, two folds do not either, and three folds return all of them,
\[
  \operatorname{spin}^{3} = \operatorname{id}, \qquad
  \operatorname{spin}^{1} \neq \operatorname{id}, \quad \operatorname{spin}^{2} \neq \operatorname{id} .
\]
The period is three, not one and not two, and the proof leans on no axiom at all.

% [4.2.p2]
Those three folds are not an abstraction laid on from outside; they are the account's own roles.
Minus, zero, and plus are the electron, the null, and the positron --- the three parts the calibration rotates through,
the same triple the whole device counts by.
The turn the machine takes on itself \emph{is} this fold, so it inherits the fold's length.
Weighing the weighing takes one step around a loop whose home is three steps away.

\section{Why three, and not one or two}

% [4.3.p1]
The count comes to rest at three because three is where the loop closes, and the two verbs say why sharply.
Each fold \emph{retanges}: it tanges the previous tange, names the naming one turn further, and selects a position the
last turn had not reached.
The first retange reaches a genuinely new place; so does the second; the third lands back on the start.
Three retanges exhaust the distinct positions the loop holds --- there are exactly three, and after them a fourth fold
only repeats the first.
A repeat carries no new characteristic, so a fourth turn does not retange; it \emph{refunges}, pooling back a place
already counted.
The count-to-three is precisely where the reflexive retange runs out: below it there is nothing new to select, so the
machine stops retangeing and lets the loop close.

\section{The residue is the loop's invariant}

% [4.4.p1]
This closing is why the leftover is the same every time, and that reproducibility is the deepest thing here.
The three costs, carried around the loop as a cycle --- the orbit cost to the pair cost, the pair to the echo, the echo
back to the orbit --- hand a residue across each edge, and the build records it.
The build proves the three edge-residues sum to zero, with no axiom,
\[
  d_1 + d_2 + d_3 = 0 .
\]
The loop \emph{funges} its three costs into one conserved quantity: it pools them around the cycle, and what the pooling
returns is unchanged from what it took in.
A quantity carried around a loop and handed back unchanged is the loop's invariant, and this leftover is exactly that:
the fixed point of a cycle whose length is provably three.

% [4.4.p2]
So the leftover does not drift from one build to the next, and now the reason is plain.
It is not a magnitude the machine happens to land on; it is the invariant a closed three-cycle conserves, and a
conserved invariant comes back the same however many times you go around.
The reproducibility that looked like luck is a theorem about a loop of length three.
And the shape is one the machine has met before at the level of the value --- a count that halts at three and refuses to
refine past its own floor --- here lifted to the level of the self-account: the same closure, read reflexively.

% [4.4.p3]
One guard belongs here, because a later step pays the coupling and this step must not seem to.
What closes at three here is a \emph{structure}, not a number.
The account settles; the loop shuts; the invariant is conserved --- and none of that is the value of the coupling, which
the machine reads as a bracket only later, and does not touch here.
The machine can be made to hand back a self-crossing ratio in passing, a small machine-dependent figure, and it must not
be mistaken for the coupling's own crossing, which is a different quantity, read elsewhere.
The two stay apart: here the loop closes at three; the number the loop is eventually asked for waits, and is named only
where the construction has earned it.

% [4.4.p4]
Derived and read, kept apart as always.
That the fold has period exactly three, that the three edge-residues sum to zero, that the account settles to a
three-count fixpoint --- these are theorems off the build.
That the self-reading \emph{is} this fold, that its closing is the loop coming home, that the three roles are electron,
null, and positron --- that is the gloss laid over the theorems, marked as such.
The account of the account does not run away.
It goes around three times, comes home, and hands back the same leftover it started with --- the price of asking,
counted once and conserved.
